\project{16}{Where the classifier runs: sensor, MCU or host}{NUCLEO-H7A3ZI-Q with X-NUCLEO-IKS4A1 and a Raspberry Pi}{In-sensor engine against on-MCU model against host, measured}

\begin{keyfacts}
\item[Board] NUCLEO-H7A3ZI-Q with X-NUCLEO-IKS4A1, and a Raspberry Pi as the host. One shield only
\item[Peripherals] I2C1 to the shield, the FIFO watermark line, USART3 as the framed link to the host, TIM2 and the cycle counter as instruments, one marker pin
\item[Toolchain] arm-none-eabi-gcc with CMake, a permissive decision-tree generator, the signal-processing library, and the in-sensor toolchain for the third leg
\item[Operating system] Bare metal
\item[Difficulty] 5 of 5
\item[Effort] 6 evenings of about four hours
\item[Deliverable] One labelled dataset, one feature definition, one trained model, and the same classification running in three places with latency, memory and charge per decision measured for each
\end{keyfacts}

\subsection*{Why this project}

Every discussion of intelligence on small systems eventually reaches the
question of where the decision should be computed, and almost every answer is an
assertion. The three candidate placements are the sensor, which now contains a
programmable core of its own; the microcontroller, which has 1312 kB of declared
static memory and a signal-processing instruction set; and a host at the other
end of a link. Each
placement has an obvious argument in its favour and no published measurement
that compares all three on one task, on one set of hardware, with a stated
method. That gap is the reason this chapter exists.

The method is deliberately narrow so that the comparison means something. One
labelled dataset, collected once. One feature definition, written once. One
model, trained once. Then the same model is compiled three ways and placed in
three locations, and three quantities are measured for each: latency from the
end of a data window to a usable class, memory occupied, and charge consumed per
decision. Anything that differs between the legs other than the placement is a
confound, and the design of this chapter is mostly the work of removing
confounds.

There is one recent preprint that measures two of the three placements and
reaches a counter-intuitive conclusion: the host wins, because the cost is
dominated by the bus read rather than by the classification. Reading a register
over this bus is three transactions, not one, and if the decision is made in the
sensor the microcontroller still has to wake up and perform them. That is a
genuine insight and it is worth testing. It is also unreviewed, its authors are
unaffiliated on the page, and its host is a far heavier machine than anything on
this bench. Treat it as the hypothesis, not the answer. A host that is itself a
microcontroller, with a wake-up measured in microseconds rather than in tens of
milliseconds, could plausibly reverse the ranking, and naming that is part of
this chapter's contribution.

\begin{note}{What this chapter has and has not produced}
The result does not exist yet. What follows is the method that produces it, and
every measured column in every table in this chapter stays blank until an
instrument has filled it. That is the standing rule of this book and it applies
with particular force here, because the whole value of the chapter is that its
numbers are new. A chapter that filled them in from plausible expectation would
be worth less than no chapter at all.
\end{note}

\subsection*{Prior art and what to reuse}

\begin{table}[H]\centering\small
\begin{tabular}{p{34mm}p{46mm}p{38mm}p{20mm}}
\toprule
Source & What it gives & What it does not & Licence \\
\midrule
The in-sensor processor examples & Working programs for the programmable core in the shield's inertial part, including prebuilt outputs so a demonstration runs before the toolchain is installed & It declares no licence at its root, so files are checked one at a time and nothing is vendored on faith & None at root \\
The ready-made classifier configurations for the in-sensor engine & Complete worked configurations for the other in-sensor engine, and the clearest statement of what that engine can express & It replaced a predecessor that was withdrawn in 2025, so a great deal of existing writing points at addresses that no longer resolve & BSD-3-Clause \\
The in-sensor state machine collection & The same for the state-machine engine, which is the right tool when the decision is a sequence rather than a classification & Nothing about features or training & BSD-3-Clause \\
A permissive decision-tree and anomaly generator for small targets & C output for trees, forests and naive Bayes from about 2 kB of flash, which is the model this chapter compiles three ways & No feature extraction and no data collection & MIT \\
The signal-processing library & The transforms and statistics the feature stage is built from, with a Python wrapper whose API mirrors the C API so both sides run the same test vectors & No model and no training & Apache-2.0 \\
The permissive neural runtime & A clean path for the neural comparison leg, with optimised kernels that do help on this core through its signal-processing instructions & It is the comparison, not the baseline, and its speed-up figures in vendor material belong to later cores & Apache-2.0 \\
One recent preprint on placement & The hypothesis, and a measurement of two of the three placements & Review, affiliation, and a host anywhere near this size & Reading only \\
\bottomrule
\end{tabular}
\caption{Prior art for chapter 16. The repository with no licence at its root is the one most likely to end up copied by accident, because it is the one that ships ready-to-run outputs.}
\end{table}

What is left to write is the comparison itself, which is the part nobody has
published, plus the two pieces of discipline that make it a comparison rather
than three demonstrations: a feature definition that produces bit-identical
values in all three implementations, and a measurement protocol that states what
is inside and outside each boundary.

\begin{note}{Writing that points at addresses which no longer resolve}
Both in-sensor repositories replaced predecessors that were discontinued during
2025. A large amount of tutorial material, including material that is otherwise
good, still links to the old locations. When a search result for this engine
looks authoritative and its links do not resolve, that is the reason, and the
current repositories are named in the sources at the end of this chapter.
\end{note}

\subsection*{Parts from the inventory}

\begin{table}[H]\centering\small
\begin{tabular}{p{40mm}p{66mm}p{40mm}}
\toprule
Part & Role & Interface \\
\midrule
NUCLEO-H7A3ZI-Q & The middle placement, and the bus master and link driver for the other two & Micro USB to the host \\
X-NUCLEO-IKS4A1 & The sensor, including the inertial part with the programmable core that hosts the first placement. The only shield on the board & I2C1 plus the watermark line \\
Raspberry Pi 4 & The third placement, at the far end of a framed serial link & USB serial, or the header serial pins \\
Renkforce USB/TTL cable & The link between the board and the host when the header pins are used. The red 5 V lead is disconnected & 3.3 V serial \\
nRF-PPK2 & Charge per decision for the node. It measures up to 1 A and the host is outside its range and outside this measurement & IDD jumper and one digital input \\
Host PC & Data collection, training, and the build for all three legs & Windows with the Linux subsystem \\
\bottomrule
\end{tabular}
\caption{Inventory items used in chapter 16. Nothing is bought. The single-board computer is the host under test and is deliberately not the machine that trains the model, so that the training environment is not accidentally part of the comparison.}
\end{table}

\subsection*{System architecture}

\diagram{c16_arch}{The same model in three places, and what crosses each boundary. The boundaries are where the cost is, which is the preprint's insight and the thing this chapter is built to measure.}

Read that figure by the arrows rather than by the boxes. In the first leg a
class index crosses the bus. In the second a window of samples crosses the bus
and nothing crosses the link. In the third a window of samples crosses the bus
and then crosses the link as well. The work performed is nearly identical in all
three; what differs is how much data moves and how many times something has to
wake up to move it.

\subsection*{Peripheral configuration}

\begin{table}[H]\centering\small
\begin{tabular}{p{22mm}p{26mm}p{24mm}p{30mm}p{30mm}}
\toprule
Peripheral & Mode & Clock source & Pins and function & Interrupt and transfers \\
\midrule
I2C1 & Fast mode, 400 kHz & Peripheral bus, source selected in RCC & D14 and D15 on the Arduino header. Confirm in MB1363 & Interrupt-driven. DMA is a variant and Chapter 4 owns it \\
GPIO input & Rising edge, external interrupt & Peripheral bus & FIFO watermark line, or the in-sensor result line in the first leg & EXTI, pre-emption priority 6 \\
USART3 & Asynchronous, 115200 8N1 & Peripheral bus & The framed link to the single-board computer & Polled here, DMA in Chapter 4 \\
TIM2 & Free running, 32-bit, 1 MHz & Peripheral bus & None & The timestamp on every window \\
DWT cycle counter & Free running & Core clock & None & The instrument for the compute stages \\
GPIO output & Marker pin, one pulse per stage & Peripheral bus & One free Zio pin into a PPK2 digital input & None \\
\bottomrule
\end{tabular}
\caption{Peripheral configuration. The marker pin is the instrument that makes the three legs comparable: it separates wake, read, compute and send inside one current trace, which is the method Chapter 10 builds and this chapter applies.}
\end{table}

\subsection*{Wiring}

\diagram{c16_wiring}{The shield on the board, the instrument across the board's own supply, and the host at the end of a serial link. What the instrument can and cannot see is drawn deliberately, because the boundary of the measurement is part of the result.}

The safety points are the ones this bench always has, plus one that is specific
here. Everything is 3.3 V logic. The red 5 V lead of the serial cable is
disconnected, because the board and the single-board computer both have their
own supplies and joining them is how a ground loop or a back-powered board
happens. The instrument measures up to 1 A, which is ample for the board and
shield and nowhere near enough for the single-board computer, so the host's own
energy is not measured in this chapter and is named as excluded rather than
quietly omitted.

\subsection*{Memory and timing budget}

\diagram{c16_mem}{Footprint in the three placements, drawn to make one point: on this part memory is not the binding constraint. The reference anomaly model from the standard suite is about 270 kB and fits into this part several times over. In the sensor it does not fit at all.}

\begin{table}[H]\centering\small
\begin{tabular}{p{44mm}p{28mm}p{28mm}p{30mm}}
\toprule
Quantity & Budget & Measured & Margin \\
\midrule
Feature stage, flash on the MCU & 6 kB & not measured & not measured \\
Decision tree, flash on the MCU & 4 kB & not measured & not measured \\
Neural comparison model, flash & 64 kB & not measured & not measured \\
Window buffer, SRAM & 8 kB & not measured & not measured \\
In-sensor program memory & read from the application note & not measured & not measured \\
Latency, sensor placement & 5 ms & not measured & not measured \\
Latency, MCU placement & 5 ms & not measured & not measured \\
Latency, host placement & 30 ms & not measured & not measured \\
Charge per decision, node & not budgeted & not measured & not measured \\
Charge per decision, host & outside the instrument & not measured & excluded by design \\
\bottomrule
\end{tabular}
\caption{The budget table. The last row is the honest statement of what this bench cannot do, and it is also the row that stops this chapter from claiming to have settled the preprint's question in the general case.}
\end{table}

\begin{note}{Memory is almost never the binding constraint on this part}
This part declares 1312 kB of static memory across five regions, plus a tightly
coupled instruction memory that neither of the machine-readable sources for this
die states, and 2048 kB of flash. The commonly repeated figure of about 1.4
megabytes is reached only by adding that instruction memory in, which is why
this book states the declared total and names the addition separately. Either
way the part sits at the large end of the envelope that small-system machine
learning literature assumes. Most of that literature was written for parts an order of
magnitude smaller, and its central anxiety, fitting the model, does not apply
here. What binds instead is throughput and energy: one core, no accelerator, and
no vector extension. This core does have the signal-processing instruction
extension, so the optimised neural kernels genuinely help, through that path
rather than through vectors. Every impressive vendor speed-up figure for those
kernels belongs to a later core than this one, and a chapter that implied
otherwise would be repeating the error this whole volume is about.
\end{note}

\subsection*{Firmware design (UML)}

\diagram{c16_uml}{One feature definition, one trained model, three implementations. The dashed arrows are claims of conformance; the gate underneath is what turns a claim into evidence, and it compares outputs rather than source.}

The design is a compiler problem more than a firmware problem. The feature
definition exists once, as a specification and as a set of test vectors. Three
implementations claim to satisfy it: one in the in-sensor toolchain, one in C
against the signal-processing library, and one in Python on the host. The tree
exists once as trained coefficients and is emitted three times. Before any
measurement is taken, all three implementations are fed the same recorded window
and must produce the same class, on every window in the test set. Until that
holds, the three legs are three different experiments and their numbers cannot
be compared.

That equivalence is easier to state than to obtain, and the usual cause of
failure is numeric rather than logical. The in-sensor leg may work in a fixed
point format the other two do not use. The host may compute in double precision
where the microcontroller computes in single. The fix is to define the feature
values in the format the most constrained implementation can produce and to
require the other two to match it, rather than the other way around, which is
the opposite of what is comfortable and the only version that works.

\subsection*{Data flow (ASCII)}

\begin{asciiart}
  collection, once               training, once                 deployment, three ways
  +------------------------+     +------------------------+     +------------------------+
  | shield to board        |     | features in Python,    |     | A: features and tree   |
  |   raw labelled windows +---->|   mirroring the C API  |     |    in the sensor core  |
  |   stored on the host   |     | tree fitted once,      +---->| B: features and tree   |
  |                        |     | exported three ways    |     |    on the MCU          |
  |                        |     |                        |     | C: window over the     |
  +------------------------+     +------------------------+     |    link, tree on host  |
                                                                +------------------------+
                                              |
                                              v
  +--------------------------------------------------------------------------------------+
  | equivalence gate: all three produce the same features and the                        |
  |   same class on every window of the held-out set                                     |
  +--------------------------------------------------------------------------------------+
                                              |
                                              v
  +--------------------------------------------------------------------------------------+
  | measurement: DWT cycles, TIM2 timestamps, marker pin into the                        |
  |   instrument. Latency as median, p99 and maximum; flash and                          |
  |   SRAM; charge per decision. The host's own energy is outside                        |
  |   the instrument and is named as outside the comparison.                             |
  +--------------------------------------------------------------------------------------+
\end{asciiart}

\subsection*{Repository layout}

\begin{plaincode}
nucleo-h7a3-classifier-placement/
  CMakeLists.txt
  dataset/
    raw/                        # labelled windows, one file per session
    protocol.md                 # how each class was produced, in detail
  features/
    features.md                 # the definition, and the fixed-point format
    features.c  features.h      # the C implementation, CMSIS-DSP
    features.py                 # the Python implementation, same test vectors
    vectors/                    # inputs and expected outputs, shared by all
  model/
    train.py                    # fit once, export three ways
    tree_mcu.c                  # generated, permissive generator
    tree_host.py
    ispu/                       # the in-sensor program and its build
  src/leg_a_sensor.c            # read a class index, nothing else
  src/leg_b_mcu.c               # window, features, tree, on the board
  src/leg_c_host.c              # window, frame, send, wait for the class
  src/measure.c                 # DWT, TIM2, marker pin
  tools/equivalence.py          # the gate: three implementations, one answer
  tools/collect.py  tools/report.py
  docs/PROVENANCE.md            # every dependency, its licence, the date
  README.md
\end{plaincode}

\subsection*{Steps}

\begin{steps}
\item \textbf{Choose a task the bench can actually produce and write the
protocol down.} Four classes of movement of the board and shield on a table:
still, sliding, tapped, and lifted and rotated. Write
\texttt{dataset/protocol.md} describing exactly how each class is produced,
including how long each session runs and who produced it. Hand-generated data
limits what can be claimed about accuracy, and the protocol is what makes the
limitation visible rather than hidden. The accuracy of the classifier is not
this chapter's result; the placement comparison is.

\item \textbf{Collect the dataset once, and never collect it again.} Stream raw
windows over the link with timestamps and labels, and store them on the host.
Every later leg is fed from these files, including the legs that run on the
board, because a leg fed from live movement is being fed different data from the
others and the comparison is void.
\begin{shellcode}
python tools/collect.py --port /dev/ttyACM0 --label still --seconds 120
python tools/collect.py --port /dev/ttyACM0 --label sliding --seconds 120
python tools/collect.py --port /dev/ttyACM0 --label tapped --seconds 120
python tools/collect.py --port /dev/ttyACM0 --label rotated --seconds 120
\end{shellcode}

\item \textbf{Define the features in the most constrained format first.} Decide
the window length, the overlap and the feature list, and decide the numeric
format that the in-sensor implementation can produce. Then write that down as
the definition and make the other two implementations match it.
\begin{ccode}
/* One window: 128 samples of three axes at 104 Hz, about 1.23 s.
   Features are computed in Q15 because that is what the most
   constrained of the three implementations can produce. */
#define WIN_N      128u
#define WIN_AXES     3u
typedef struct {
    q15_t mean[WIN_AXES];
    q15_t rms[WIN_AXES];
    q15_t p2p[WIN_AXES];          /* peak to peak */
    q15_t band[WIN_AXES][4];      /* energy in four bands, from the FFT */
} features_t;
\end{ccode}

\item \textbf{Build the equivalence gate before building any leg.} This is the
step that turns three demonstrations into one experiment, and doing it later
never works because by then each implementation has its own quiet assumption.
\begin{pycode}
# Every implementation reads the same vectors and must agree exactly.
for name, impl in (("c", run_c_features), ("py", run_py_features),
                   ("ispu", run_ispu_features)):
    for vec in load_vectors("features/vectors"):
        assert impl(vec.input) == vec.expected, (name, vec.id)
\end{pycode}

\item \textbf{Train once and export three times.} The tree is fitted on the host
with the Python feature implementation, and the same fitted coefficients are
emitted as C for the microcontroller, as a program for the in-sensor core, and
as a Python object for the host leg. A tree with a bounded depth is chosen
deliberately, because bounded depth is bounded latency, and bounded latency is
half of what makes a classifier acceptable in a control loop.

\item \textbf{Leg A, in the sensor.} Load the in-sensor program, configure the
part to run it on its own data, and have the microcontroller do nothing but wait
for the line and read the result. Start from the prebuilt outputs in the
examples repository, which run before the toolchain is installed and prove the
path works, then replace them with your own. Audit the licence of every file you
keep, because the repository declares none at its root.
\begin{ccode}
void EXTI_result_handler(void)          /* leg A: the whole of the work */
{
    marker_pulse(MARKER_WAKE);
    const uint32_t t = timebase_ticks();
    uint8_t klass;
    i2c1_read_reg(IMU_ADDR, ISPU_RESULT_REG, &klass, 1);  /* three transactions */
    decision_post(t, klass);
    marker_pulse(MARKER_DONE);
}
\end{ccode}
Note what the comment records. A one-byte register read is an address write, a
repeated start and a data read. That is the cost the preprint says dominates,
and here it is, in the leg that was supposed to avoid work.

\item \textbf{Leg B, on the microcontroller.} Drain the FIFO on the watermark
line into a window, compute the features with the signal-processing library, and
evaluate the generated tree. Measure the two compute stages separately, because
the feature stage and the model are different sizes of problem and a single
number hides which one matters.
\begin{ccode}
const uint32_t c0 = DWT->CYCCNT;
features_t f;  features_compute(&win, &f);
const uint32_t c1 = DWT->CYCCNT;
const uint8_t k = tree_predict((const q15_t *) &f);
const uint32_t c2 = DWT->CYCCNT;
printf("feat %lu cyc, tree %lu cyc\r\n",
       (unsigned long)(c1 - c0), (unsigned long)(c2 - c1));
\end{ccode}

\item \textbf{Leg C, on the host.} Send the window over the framed link from
Chapter 5 and wait for the class to come back. Measure the round trip from the
board's side, which is the only side that matters for the node's latency, and
record the host's own processing time separately so that the two are not
confused.
\begin{pycode}
# On the single-board computer. Same features, same tree, different machine.
win = codec.decode(frame)
f = features.compute(win)          # the Python implementation, gated above
k = tree.predict(f)
port.write(codec.encode_class(k))
\end{pycode}

\item \textbf{Add the neural comparison, and hold it to the same gate.} Train a
small network on the same features, run it on the microcontroller under the
permissive runtime with the optimised kernels enabled, and compare it against
the tree on three axes: accuracy, kilobytes, and the spread of its latency. The
thesis worth testing is that the transform plus the statistical model wins on
accuracy per kilobyte, is deterministic in latency, and can be explained to a
customer. One peer-reviewed paper argues exactly that, and this is the chapter
that can check it on this part.

\item \textbf{Measure all three legs with one protocol and publish the
protocol.} For each leg: one hundred decisions, the marker pin pulsed at wake
and at done, the current trace captured, and the cycle counts recorded. Report
median, ninety-ninth percentile and maximum rather than a mean, because the
maximum is what a control loop has to tolerate. State which boundaries are
inside each measurement, and state that the host's own energy is outside all of
them.
\begin{shellcode}
python tools/report.py --leg a --n 100 --out results/leg_a.json
python tools/report.py --leg b --n 100 --out results/leg_b.json
python tools/report.py --leg c --n 100 --out results/leg_c.json
python tools/report.py --compare results/ --out results/README.md
\end{shellcode}
\end{steps}

\subsection*{Build, flash and debug}

\diagram{c16_timing}{One decision in each of the three placements, on one time axis. The shaded stages are what the marker pin separates inside the current trace. Every interval is a budget until an instrument has replaced it.}

\begin{shellcode}
cmake -B build-fw -G Ninja -DCMAKE_TOOLCHAIN_FILE=cmake/arm-none-eabi.cmake -DLEG=B
cmake --build build-fw -j
cp build-fw/placement.bin "$PROBE_DISK"/   # onto the probe disk
python tools/equivalence.py        # must pass before any result is recorded
\end{shellcode}

Recovering a board that will not answer is Chapter 1's material. The failure
specific to this chapter is subtler and does not present as a failure at all:
three legs that produce slightly different classes and a comparison that is
quietly meaningless. The equivalence gate is the only defence, and it belongs in
the build rather than in a habit.

\begin{note}{When the three legs disagree}
In order of likelihood: the window boundaries differ, because one leg starts its
window at a FIFO watermark and another at a sample count; the numeric format
differs, and a rounding rule that is invisible in one implementation changes a
feature at the last bit; the sample rate is not what was configured, because the
sensor's actual rate and its nominal rate differ by a few tenths of a percent;
one implementation normalises and another does not; and the tree was exported
from a different fit than the one deployed. Check them by feeding a recorded
window through all three and comparing the features before the class, which
localises the disagreement to a stage instead of to a leg.
\end{note}

\subsection*{Verification and acceptance criteria}

\begin{itemize}
\item The equivalence gate passes: all three implementations produce identical
features and identical classes on every window of the held-out set. The build
fails if it does not.
\item Every number in the results directory names the instrument that produced
it, and every number that has not been produced is absent rather than estimated.
\item Latency is reported as median, ninety-ninth percentile and maximum over at
least one hundred decisions per leg, never as a mean.
\item Memory is reported from the size output for the microcontroller legs and
from the in-sensor toolchain's own output for the sensor leg, with the two
clearly not being the same kind of number.
\item Charge per decision for the node is reported for all three legs, and the
exclusion of the host's own energy is stated in the same table rather than in a
footnote elsewhere.
\item The neural comparison is reported on accuracy, kilobytes and latency
spread, and the conclusion is stated even if it contradicts the thesis this
chapter set out to test.
\item The dataset protocol is published with the dataset, so that a reader can
see exactly how limited the hand-generated data is.
\end{itemize}

\subsection*{Variants}

\begin{table}[H]\centering\small
\begin{tabular}{p{24mm}p{26mm}p{44mm}p{22mm}p{20mm}}
\toprule
Axis & Variant & What changes & Cost & Built in full in \\
\midrule
Intelligence and reach & Classifier in the sensor & The decision is made on the sensor's own core and the microcontroller reads a class index. The bus read it still has to perform is the thing under test & A toolchain, and a repository with no licence at its root & Here \\
Intelligence and reach & Model on the MCU & Raw windows cross the bus, the features and the tree run on the board, and nothing crosses the link & Core time and flash, neither of which is scarce here & Here \\
Intelligence and reach & Model on the host & The window crosses the link and the class comes back. The node becomes a sensor and a modem & Link energy, and a host that must be awake & Here. The framed link itself is Chapter 9 \\
Intelligence and reach & Local only & No host at all, which is the baseline the first two legs already are & Nothing & Chapters 8 and 13 \\
Execution model & DMA into the window & The FIFO drain stops occupying the core, which matters most in the leg that has the most work to do & Cache maintenance on the window buffer & Chapter 4 and Chapter 19 \\
Operating system & An RTOS task set & The three stages become tasks with stated priorities and a measured latency distribution, which is a better home for the ninety-ninth percentile question & A scheduler to justify & Chapter 20 \\
Language & Host in Python & The host leg as written. A compiled host would change the host's numbers and not the node's & Host-side latency & Chapter 3 \\
\bottomrule
\end{tabular}
\caption{Variants for chapter 16. The three placements are this chapter's own, which is why the whole comparison lives here rather than being spread over three chapters where nothing could be compared.}
\end{table}

\subsection*{Pitfalls}

\begin{itemize}
\item Comparing three legs that were fed different data. A leg driven by live
movement is not comparable with a leg driven by a recording, however carefully
the movement is repeated.
\item Letting the numeric formats differ. The equivalence has to be defined in
the format the most constrained implementation can produce, which is not the
format that is most convenient to train in.
\item Reporting a mean latency. The interesting quantity is the tail, and the
mean hides exactly the behaviour that decides whether a classifier can sit
inside a control loop.
\item Assuming the in-sensor leg does no work on the microcontroller. It still
wakes, still performs a multi-transaction bus read, and that cost is the
preprint's whole argument.
\item Vendoring from a repository that declares no licence at its root because
its prebuilt outputs made the demonstration easy. Read it, run it, and check
each file before keeping any of it.
\item Quoting a vendor speed-up figure for the optimised neural kernels. Those
figures belong to later cores with a vector extension this part does not have.
\item Treating the popular training service as a supported path. It treats this
part as a porting target rather than a supported board: no ready-made firmware,
no data forwarder, and an exported kit that declares no single machine-readable
licence, with one optimising component proprietary and another sold only at the
top tier. The permissive runtime is the clean path and this chapter uses it.
\item Presenting the host's energy as unmeasured when it is in fact
unmeasurable on this bench. Those are different statements and the second one
has to be said out loud.
\item Claiming the comparison settles the question in general. It settles it for
this task, this data, this model and these three machines, which is already more
than anyone has published.
\end{itemize}

\subsection*{Best practices applied}

\begin{itemize}
\item The experiment is designed to remove confounds before it is designed to
produce numbers, and the equivalence gate is in the build rather than in the
author's discipline.
\item Every dependency's licence is stated before code is written around it, and
the one dependency with no licence at its root is named as such in the same
sentence that recommends reading it.
\item The boundary of the measurement is published with the measurement. What
the instrument cannot see is part of the result.
\item The claim is scoped to what was measured, and the configuration that could
overturn it, a microcontroller host with a short wake-up, is named rather than
left for a reader to discover.
\item Tail statistics rather than means, which is the reporting method the
serious real-time measurement literature uses and almost no vendor material
does.
\end{itemize}

\subsection*{Stretch goals}

\begin{itemize}
\item Add a fourth leg with a microcontroller as the host, using one of the
radio coprocessor boards on this bench over a serial link. Its wake-up is orders
of magnitude shorter than the single-board computer's, and it is the
configuration most likely to reverse the preprint's ranking. That result would
be new and it is within reach of this bench.
\item Run the standard small-system machine learning benchmark suite on this
part. No published result exists for this core: the suite's reference baseline
is a smaller core, the recent vendor submissions are a different core and an
accelerated preview part, and the vendor's own harness for running it has been
dormant since 2024.
\item Repeat the comparison with the state-machine engine rather than the
classifier engine for the in-sensor leg, which suits a decision that is a
sequence rather than a snapshot and would show whether the placement ranking
depends on the kind of decision.
\item Replace the hand-generated dataset with a mechanically repeatable one,
which is the single change that would turn the accuracy numbers from
illustrative into citable.
\end{itemize}

\subsection*{Roadmap and next steps}

Chapter 17 changes the subject from where code runs to how much a layer of
abstraction costs, and it uses the same instruments: the cycle counter, the size
output and the marker pin. Chapter 18 builds the transform stage this chapter
depends on, with a numerical acceptance test, and Chapter 20 gives the three
stages somewhere sensible to live when they have to share a processor with
anything else.

For going deeper on the machine-learning side, the free open textbook that grew
out of a university course and an online certificate is the closest thing to a
canonical curriculum for this material, and it is the right next step after this
chapter rather than before it. The free-to-audit course on embedded machine
learning, taught by the author of several widely used video series, covers the
same ground with exercises. The community embedded engineering roadmap is the
map for everything around it, and its central claim, that projects outrank
reading, is the reason this chapter is a build rather than a survey. Appendix H
lists them with their licences, including the two whose terms make them
something to recommend and never to quote.

\subsection*{Portfolio evidence}

\begin{itemize}
\item A public repository containing a labelled dataset, its collection
protocol, one feature definition with shared test vectors, and three
implementations that provably agree.
\item A results directory with latency, memory and charge for three placements,
each number naming its instrument, and a stated boundary for what the instrument
could not see.
\item A short written result that says which placement won on this task and by
how much, what the preprint predicted, and where the two agree and disagree.
That is a new result and it should be presented as one.
\item The equivalence gate running in the build, which is the piece of
engineering judgement in this chapter that a reviewer will recognise fastest.
\end{itemize}

\subsection*{Sources}

Normative references:
\begin{itemize}
\item UM3239, the user manual for the X-NUCLEO-IKS4A1, for the sensor
complement, the addresses and the three bus topologies.
\item The datasheet and the application note for the inertial part that carries
the programmable core, for the core's program memory size, its instruction set
and how a program is loaded.
\item Reference manual RM0455 for I2C1, USART3, the cycle counter and the timer
on this part. Not RM0433, which documents its better-known sibling.
\item The board user manual for MB1363, for the Arduino header mapping and the
free Zio pins that carry the marker signals.
\end{itemize}

Reusable implementations:
\begin{itemize}
\item The in-sensor processor examples, which ship prebuilt outputs so a
demonstration runs before the toolchain is installed. No licence is declared at
the repository root, so every file is checked individually.\newline
\url{https://github.com/STMicroelectronics/st-mems-ispu}
\item Ready-made configurations for the in-sensor classifier engine,
BSD-3-Clause. This repository replaced one that was discontinued in 2025.\newline
\url{https://github.com/STMicroelectronics/st-mems-machine-learning-core}
\item The in-sensor state machine collection, BSD-3-Clause, with the same
history of replacing a discontinued predecessor.\newline
\url{https://github.com/STMicroelectronics/st-mems-finite-state-machine}
\item The permissive generator for trees, forests, naive Bayes and anomaly
detection on small targets, MIT, from about 2 kB of flash.\newline
\url{https://github.com/emlearn/emlearn}
\item The signal-processing library, Apache-2.0, whose Python wrapper mirrors the
C API so that both sides run the same test vectors.\newline
\url{https://github.com/ARM-software/CMSIS-DSP}
\item The optimised neural kernels, Apache-2.0, which help on this core through
its signal-processing instructions rather than through a vector extension it
does not have.\newline
\url{https://github.com/ARM-software/CMSIS-NN}
\item The permissive neural runtime, Apache-2.0, which is the clean path for the
comparison leg.\newline
\url{https://github.com/tensorflow/tflite-micro}
\item The peer-reviewed argument that classical models plus good features
outperform neural networks per kilobyte on small targets.\newline
\url{https://arxiv.org/pdf/2107.09448}
\item The unreviewed preprint that measures two of the three placements and
reaches the counter-intuitive result this chapter treats as its hypothesis.\newline
\url{https://arxiv.org/abs/2606.00524}
\item The free open textbook that is the closest thing to a canonical curriculum
for this material.\newline
\url{https://mlsysbook.ai/}
\end{itemize}
